% ═══════════════════════════════════════════════════════════════
% AB-01b: Mechanische Energie
% Physik Klasse 10 MR
% ═══════════════════════════════════════════════════════════════

\documentclass[11pt, a4paper]{article}

\usepackage[utf8]{inputenc}
\usepackage[T1]{fontenc}
\usepackage[ngerman]{babel}

\usepackage{helvet}
\renewcommand{\familydefault}{\sfdefault}

\usepackage{microtype}
\usepackage[margin=1.5cm, top=2cm, bottom=1.5cm]{geometry}
\usepackage{enumitem}
\usepackage{tabularx}
\usepackage{booktabs}
\usepackage{array}
\usepackage{tikz}
\usetikzlibrary{calc, arrows.meta, shapes.geometric}

\usepackage{amsmath, amssymb}
\usepackage{siunitx}
\sisetup{locale = DE, per-mode = symbol}

\usepackage{fancyhdr}
\usepackage{lastpage}
\usepackage{needspace}

\usepackage{xcolor}
\usepackage{tcolorbox}
\tcbuselibrary{skins}

% ═══════════════════════════════════════════════════════════════
% FARBEN
% ═══════════════════════════════════════════════════════════════

\definecolor{physik}{HTML}{2c5aa0}
\definecolor{hellgrau}{HTML}{F5F5F5}
\definecolor{mittelgrau}{HTML}{888888}
\definecolor{rahmen}{HTML}{CCCCCC}
\definecolor{Epot}{HTML}{2c5aa0}
\definecolor{Ekin}{HTML}{c62828}

\colorlet{fachfarbe}{physik}

% ═══════════════════════════════════════════════════════════════
% VARIABLEN
% ═══════════════════════════════════════════════════════════════

\newcommand{\fach}{Physik}
\newcommand{\thema}{Mechanische Energie}
\newcommand{\klasse}{10 MR}
\newcommand{\stunde}{01b}
\newcommand{\maxpunkte}{24}

% ═══════════════════════════════════════════════════════════════
% KOPF- UND FUSSZEILE
% ═══════════════════════════════════════════════════════════════

\pagestyle{fancy}
\fancyhf{}
\renewcommand{\headrulewidth}{0.4pt}
\renewcommand{\footrulewidth}{0pt}
\lhead{\small \fach{} Klasse \klasse{} | \thema}
\rhead{\small Name: \underline{\hspace{3.5cm}}}
\cfoot{\small\textcolor{mittelgrau}{Seite \thepage{} von \pageref{LastPage}}}

% ═══════════════════════════════════════════════════════════════
% BOXEN
% ═══════════════════════════════════════════════════════════════

\newtcolorbox{merkkasten}[1][]{
    colback=hellgrau,
    colframe=hellgrau,
    coltitle=black,
    colbacktitle=hellgrau,
    boxrule=0pt,
    arc=0pt,
    left=8pt, right=8pt, top=8pt, bottom=6pt,
    borderline north={2pt}{0pt}{fachfarbe},
    before skip=6pt, after skip=6pt,
    fonttitle=\bfseries,
    #1
}

\newtcolorbox{formelkasten}{
    colback=white,
    colframe=fachfarbe,
    boxrule=1.5pt,
    arc=0pt,
    left=8pt, right=8pt, top=6pt, bottom=6pt,
    before skip=4pt, after skip=4pt
}

% ═══════════════════════════════════════════════════════════════
% BEFEHLE
% ═══════════════════════════════════════════════════════════════

\newcommand{\luecke}[1]{\underline{\hspace{#1}}}
\newcommand{\punktebox}[1]{\hfill\small\textcolor{mittelgrau}{[\_\_/#1]}}

\newcommand{\aufgabe}[1]{%
    \needspace{4\baselineskip}%
    \vspace{10pt}%
    \noindent\textbf{\textcolor{fachfarbe}{Aufgabe #1}}%
    \vspace{4pt}%
    \nopagebreak[4]%
}

\newcommand{\operator}[1]{\textbf{\textcolor{fachfarbe}{#1}}}

\newlist{teil}{enumerate}{2}
\setlist[teil,1]{label=\alph*), leftmargin=1.8em, itemsep=3pt, parsep=0pt, topsep=3pt}

\newcommand{\antwortzeilen}[1]{%
    \par\vspace{4pt}%
    \foreach \n in {1,...,#1}{%
        \noindent\rule{\linewidth}{0.3pt}\vspace{6pt}%
    }%
}

\setlength{\parindent}{0pt}
\setlength{\parskip}{4pt}

% ═══════════════════════════════════════════════════════════════
% DOKUMENT
% ═══════════════════════════════════════════════════════════════

\begin{document}

% ═══════════════════════════════════════════════════════════════
% HEADER
% ═══════════════════════════════════════════════════════════════

\begin{center}
    {\Large\bfseries\textcolor{fachfarbe}{\thema}}\\[0.2em]
    {\small Arbeitsblatt | \fach{} Klasse \klasse{} | Stunde \stunde}
\end{center}
\vspace{-0.3em}
\hrule
\vspace{0.6em}

% ═══════════════════════════════════════════════════════════════
% TEIL 1: MERKKÄSTEN
% ═══════════════════════════════════════════════════════════════

\textbf{\large Teil 1: Merkkästen}

\begin{merkkasten}[title={\textbf{Kasten 1: Was ist Energie?}}]
\textbf{Energie} ist die Fähigkeit eines Körpers, \luecke{3cm} zu verrichten

oder \luecke{2cm} abzugeben.

\vspace{6pt}
\textbf{Einheit:} \luecke{3cm} (Abkürzung: \luecke{1cm})
\end{merkkasten}

\begin{merkkasten}[title={\textbf{Kasten 2: Lageenergie (potenzielle Energie)}}]
\begin{formelkasten}
\centering
\Large $E_\text{pot} = m \cdot g \cdot h$
\end{formelkasten}

\vspace{4pt}
\begin{tabular}{@{}ll@{}}
$E_\text{pot}$ & Lageenergie in \luecke{2cm} \\[4pt]
$m$ & Masse in \luecke{2cm} \\[4pt]
$g$ & Ortsfaktor: \textbf{\luecke{1.2cm} m/s²} (auf der Erde) \\[4pt]
$h$ & Höhe in \luecke{2cm}
\end{tabular}
\end{merkkasten}

\begin{merkkasten}[title={\textbf{Kasten 3: Bewegungsenergie (kinetische Energie)}}]
\begin{formelkasten}
\centering
\Large $E_\text{kin} = \frac{1}{2} \cdot m \cdot v^2$
\end{formelkasten}

\vspace{4pt}
\begin{tabular}{@{}ll@{}}
$E_\text{kin}$ & Bewegungsenergie in \luecke{2cm} \\[4pt]
$m$ & Masse in \luecke{2cm} \\[4pt]
$v$ & Geschwindigkeit in \luecke{2cm}
\end{tabular}

\vspace{4pt}
\textbf{Achtung:} $v$ wird \luecke{2.5cm}! Doppelte Geschwindigkeit $\rightarrow$ \luecke{2.5cm} Energie!
\end{merkkasten}

\begin{merkkasten}[title={\textbf{Kasten 4: Energieerhaltungssatz}}]
Energie kann nicht erzeugt oder vernichtet werden.

Sie kann nur von einer Form in eine andere \luecke{3cm} werden.

\vspace{6pt}
\begin{formelkasten}
\centering
\Large $E_\text{ges} = E_\text{pot} + E_\text{kin} = \text{konstant}$
\end{formelkasten}
\end{merkkasten}



% ═══════════════════════════════════════════════════════════════
% TEIL 2: AUFGABEN
% ═══════════════════════════════════════════════════════════════
\vspace{6pt}
\textbf{\large Teil 2: Aufgaben}

\aufgabe{1 – Energieformen berechnen}

Verwende $g = 10$ m/s².

\begin{teil}
\item ($\star$) Ein Wassereimer ($m = 10$ kg) steht auf einem Dach in $h = 8$ m Höhe. \punktebox{2}

\operator{Berechne} die Lageenergie.

\vspace{0.3em}
$E_\text{pot} = $ \luecke{8cm}

\vspace{0.8em}

\item ($\star\star$) Ein Ball ($m = 0{,}5$ kg) fliegt mit $v = 20$ m/s. \punktebox{3}

\operator{Berechne} die Bewegungsenergie.

\vspace{0.3em}
$E_\text{kin} = $ \luecke{8cm}
\end{teil}

\vspace{1em}

\newpage

\aufgabe{2 – Energieerhaltung beim freien Fall}

($\star\star$) Ein Stein ($m = 5$ kg) fällt aus $h = 20$ m Höhe. \punktebox{5}

\begin{teil}
\item \operator{Berechne} die Lageenergie $E_\text{pot}$ oben (vor dem Fallen).

\vspace{0.3em}
$E_\text{pot} = $ \luecke{8cm}

\vspace{0.6em}

\item \operator{Gib an}, wie groß die Bewegungsenergie $E_\text{kin}$ unten ist (kurz vor dem Aufprall).

\vspace{0.3em}
$E_\text{kin} = $ \luecke{3cm}

\vspace{0.6em}

\item \operator{Begründe} deine Antwort in einem Satz.

\antwortzeilen{2}
\end{teil}

\vspace{1em}

\aufgabe{3 – Energieflussdiagramm (realistisch)}

($\star\star$) Ein Kind rutscht eine Rutsche hinunter. In der Realität wird ein Teil der Energie durch Reibung in Wärme umgewandelt. \punktebox{6}

\begin{teil}
\item \operator{Zeichne} ein Energieflussdiagramm für diesen Vorgang.

\begin{center}
\begin{tikzpicture}
    % Startbox
    \draw[thick] (0,0) rectangle (2.2,1);
    \node at (1.1, 0.5) {\luecke{1.2cm}};

    % Pfeil
    \draw[thick, ->, >=Stealth, line width=2pt] (2.5,0.5) -- (3.8,0.5);

    % Zwei Zielboxen (E_kin + E_Wärme)
    \draw[thick] (4.1,0) rectangle (6.3,1);
    \node at (5.2, 0.5) {\luecke{1.2cm}};

    \node at (6.6, 0.5) {$+$};

    \draw[thick] (6.9,0) rectangle (9.1,1);
    \node at (8.0, 0.5) {\luecke{1.2cm}};
\end{tikzpicture}
\end{center}

\vspace{0.5em}

\item \operator{Erkläre}, was bei diesem Vorgang mit der Energie passiert.

Verwende die Begriffe: \textit{Lageenergie, Bewegungsenergie, Wärme, umgewandelt, Reibung}

\antwortzeilen{3}
\end{teil}

\vspace{0.8em}

\aufgabe{4 – Knobelaufgaben (Transfer)}

($\star\star\star$) Diese Aufgaben erfordern, dass du die Formeln „rückwärts" anwendest. \punktebox{8}

\begin{teil}
\item Ein Auto ($m = 1000$ kg) fährt mit $v = 10$ m/s auf einen Berg zu. Der Motor ist aus.

\operator{Berechne} zuerst die Bewegungsenergie $E_\text{kin}$, dann \operator{ermittle}, wie hoch das Auto kommt.

\vspace{0.3em}
$E_\text{kin} = $ \luecke{6cm}

\vspace{0.5em}
Das Auto erreicht eine Höhe von $h = $ \luecke{2cm}

\vspace{3cm}

\item Ein Stein ($m = 2$ kg) fällt aus $h = 5$ m Höhe.

\operator{Berechne} $E_\text{pot}$ oben, dann \operator{ermittle}, wie schnell der Stein unten ist.

\vspace{0.3em}
$E_\text{pot} = $ \luecke{6cm} $= E_\text{kin}$ (unten)

\vspace{0.5em}
Der Stein hat unten $v = $ \luecke{2cm}
\end{teil}

\vspace{1em}

% ═══════════════════════════════════════════════════════════════
% FOOTER
% ═══════════════════════════════════════════════════════════════



\end{document}
